\section{Conclusion}

Completion of this laboratory work has provided practical understanding of actuator
control techniques for embedded systems, including analog PWM motor control, binary
relay control, and a multi-stage command conditioning pipeline, while reinforcing the
use of FreeRTOS synchronization primitives for safe inter-task communication.

This seventh laboratory work covered the complete actuator control pipeline from serial
command reception through multi-stage command conditioning to hardware actuation and
structured reporting. The key takeaways are:

\begin{enumerate}
  \item Implementing a \textbf{serial command interface} in FreeRTOS context -- the
    command task polls the serial port every 50\,ms, parses case-insensitive commands,
    validates ranges, and signals the conditioning task via semaphore, keeping the
    command parsing decoupled from the conditioning logic.
  \item Applying a \textbf{four-stage command conditioning pipeline}:
    \begin{itemize}
      \item \textbf{Saturation} clamps the raw command to the valid actuator range
        ($0$--$100\,\%$), rejecting physically impossible targets at the input boundary.
      \item \textbf{Median filter} (window = 5) removes impulsive command noise --
        isolated wrong commands are eliminated without distorting the underlying
        command trend, as the median is inherently robust to up to
        $\lfloor n/2 \rfloor$ outliers in the window.
      \item \textbf{Exponential Moving Average} ($\alpha = 0.3$) provides weighted
        smoothing where recent commands contribute 30\% and the accumulated history
        contributes 70\%, reducing abrupt step changes visible to the actuator.
      \item \textbf{Ramping} ($2\,\%$ per 50\,ms cycle) limits the rate of change
        of \texttt{currentSpeed}, protecting the motor from sudden current surges and
        implementing smooth acceleration and deceleration ($\approx 2.5\,\mathrm{s}$
        for a full 0--100\,\% ramp). The motor state machine (IDLE / RAMP\_UP /
        RUNNING / RAMP\_DOWN) makes the ramp phase clearly visible on the OLED and
        in the serial plotter.
    \end{itemize}
  \item Implementing an \textbf{emergency stop} (\texttt{X}) that bypasses the
    conditioning pipeline entirely -- \texttt{currentSpeed} is set to zero immediately,
    all filter state (median buffer, EMA accumulator) is reset, and the motor stops
    without any ramp. This demonstrates the importance of having a safety override
    path independent of the normal control loop.
  \item Controlling an \textbf{analog actuator} (DC motor via L298N) using the ESP32
    LEDC peripheral: channel~0, 5\,kHz, 8-bit resolution, with direction pins managed
    separately. The MotorDriver API (Init / SetSpeed / SetDirection / Stop / GetDuty)
    is structured so that swapping the test LED for a real motor requires zero code
    changes.
  \item Controlling a \textbf{binary actuator} (relay) with a simple ON/OFF API and
    state tracking in shared data. The relay is only written when its target differs
    from the current state, avoiding redundant GPIO writes.
  \item Evaluating \textbf{alert conditions} (OVERLOAD $\geq 95\,\%$,
    LIMIT\_REACHED at boundaries) on the conditioned speed and driving LED indicators
    accordingly, providing clear visual feedback about the actuator's operating region.
  \item Structuring the firmware into four distinct \textbf{FreeRTOS tasks} with
    well-defined roles and priorities: Command (priority 3) ensures serial input is
    never missed, Conditioning and Actuator (priority 2) run the control loop, and
    Reporter (priority 1) provides periodic human-readable output without interfering
    with the control loop.
  \item Using a \textbf{mutex-protected shared data store} (\texttt{ActuatorData}) as
    the single source of truth for all actuator state -- target, conditioned, and
    current speed; direction; motor state; relay states; and alerts -- preventing
    data races when four tasks read or write concurrently.
  \item Organizing the project into \textbf{three hardware-abstraction layers} (MCAL,
    ECAL, SRV) with PlatformIO, keeping all hardware-specific code in the MCAL layer
    and making the service layer fully portable across different MCU platforms.
\end{enumerate}

This laboratory work demonstrates how a pipeline of classical command processing
techniques -- saturation, median filtering, exponential moving averaging, and ramping --
combined with FreeRTOS synchronization primitives, transforms raw serial commands into
smooth, safe actuator control suitable for both analog (PWM motor) and binary (relay)
actuators in an embedded system.
